\section{Main Design Choices}


\subsection{Tracking Ongoing Operations}
To handle the asynchronous nature of a distributed system, the
architecture avoids blocking threads while waiting for network responses.
Instead, ongoing operations are tracked using stateful context objects stored
in the node's local memory. 

This stateful tracking mechanism is utilized for both external Client
Read/Write requests and internal Join/Recovery synchronizations.

When a node coordinates a transaction—either on behalf of an external client or
on behalf of itself during a reboot—it generates a unique transaction
identifier and creates a temporary tracking object. This object holds essential
information: a reference to the requester, the target data key, and a dynamic
collection of the responses arriving from the replicas. 

These tracking objects are stored in in-memory data structures. For
external client operations, the Coordinator simultaneously initiates an
asynchronous timer based on a predefined timeout threshold. 

As acknowledgment messages arrive from the network, the node uses the
transaction identifier to retrieve the corresponding tracking object, updates
its internal counters, and evaluates whether the required Read or Write Quorum
has been satisfied. 

Once the quorum is reached, the transaction is finalized: if it was a client
request, the appropriate response is sent back to the user; if it was an
internal background read for a joining node, the node safely commits the
verified data to its own disk. The tracking object is then safely deleted from
memory. Furthermore, if a client operation's asynchronous timer expires before
the quorum is achieved, the operation is automatically aborted, the tracking
object is cleaned up, and the client is notified of the timeout, ensuring that
the system never hangs indefinitely.

\subsection{Main Messages Used in the System}
To keep the code simple and safe, all messages in the system are implemented as
Java \texttt{Records}. We can group them into three main categories based on
their purpose:

\begin{itemize}
    \item \textbf{Client Messages:} Used for communication between the user and
        the system.
    \begin{itemize}
        \item \texttt{ClientGetRequestMsg} / \texttt{ClientUpdateRequestMsg}:
            Sent by the client to ask the system to read or write data.
        \item \texttt{ClientGetResponseMsg} / \texttt{ClientUpdateResponseMsg}:
            Sent back to the client by the Coordinator to deliver the final
            result (success, failure, or data).
    \end{itemize}
    
    \item \textbf{Internal Storage Messages:} Used between the Coordinator and
        the Replicas to manage Quorums and locks.
    \begin{itemize}
        \item \texttt{ReplicaReadRequestMsg} / \texttt{ReplicaWriteRequestMsg}:
            The Coordinator sends these to ask replicas for data or to save new
            data.
        \item \texttt{ReplicaReadResponseMsg} /
            \texttt{ReplicaWriteResponseMsg}: The replicas reply with the
            requested data or confirm the save. The read response also tells
            the Coordinator if the key was already locked by someone else.
        \item \texttt{UnlockKeyMsg}: Used to free a locked key after a write
            finishes or a timeout happens.
    \end{itemize}
    
    \item \textbf{Ring Management Messages:} Used to update the network shape
        and move data without losing it.
    \begin{itemize}
        \item \textbf{Topology Updates:} Messages like \texttt{JoinMsg},
            \texttt{LeaveMsg}, and \texttt{JoinedNetworkMsg} are used to add or
            remove nodes and share the new map of the network.
        \item \textbf{Data Transfers:} Messages like \texttt{StorageRequestMsg}
            and \texttt{NodeLeavingMsg} are used to safely pass data from one
            node to another when the ring changes.
    \end{itemize}
\end{itemize}

\subsection{Simulating Crash and Recovery}
In a real distributed system, if a server loses power, its active memory (RAM)
is cleared, but the data on its physical hard drive is kept safe. To simulate
this correctly in our Akka framework, we must separate temporary memory from
permanent storage.

\textbf{1. Simulating a Crash}
When the Manager wants to simulate a hardware failure, it sends a
\texttt{CrashMsg} to a specific node. The node immediately executes the
\texttt{actorContext.stop(self)} command. This completely destroys the Akka
Actor. All of its temporary memory—like pending queues, locks, and network
connections—is instantly erased, perfectly simulating a sudden power loss.

\textbf{2. Simulating a Physical Hard Drive}
To ensure that the saved data survives the crash, we use a \texttt{static
ConcurrentHashMap} named \texttt{persistentDisks}. Because this map is
\texttt{static}, it exists outside the lifecycle of any single Actor. It acts
like a warehouse of physical hard drives, where each box is labeled with a Node
ID. When an Actor dies, its "hard drive" stays safely in the warehouse.

\textbf{3. The Recovery Process}
When the Manager decides to reboot the crashed node, it spawns a new Akka
Actor with the same ID and sends it a \texttt{RecoverMsg}. The recovery happens
in three steps:
\begin{itemize}
    \item \textbf{Mounting the Disk:} The new Actor checks the
        \texttt{persistentDisks} map, finds the hard drive matching its ID, and
        loads its previous data.
    \item \textbf{Updating the Topology:} The node contacts a bootstrap peer
        (another active node) to get the latest map of the network ring.
    \item \textbf{Synchronizing Missed Data:} The recovering node asks its
        clockwise neighbor for data. It then combines the neighbor's keys with
        the keys on its own hard drive. For every key it is responsible for, it
        launches a background Quorum Read. This ensures that if any data was
        updated by the network while the node was dead, the node safely
        downloads the newest version before it officially resumes normal
        operations.
\end{itemize}

